\section{Setup}

Regge calculus~\cite{Regge:1961ct} is a discrete gravitational theory on simplicial manifolds $\Delta$ with the length of edges $\{l_e\}$ dynamical. In a Lorentzian setting, a simplicial manifold consists of intrinsically flat Lorentzian $4$-simplices $\sigma\in\Delta$. Curvature is located at triangles $t\in\Delta$ and captured by deficit angles $\delta_t(\{l_e\})$, lying in the space orthogonal to $t$. For $t$ spacelike (timelike), this space is isomorphic to $\R^{1,1}$ ($\R^2$).\footnote{For the remainder, we neglect lightlike edges, triangles and tetrahedra. See~\cite{Sorkin:2019llw} for the definition of deficit angles associated to $(d-1)$-dimensional null cells.} Lorentzian deficit angles are defined as $\delta_t^{\mathrm{L},\pm} = \mp i 2\pi  -\sum_{\sigma\supset t}\psi_{\sigma,t}^{\mathrm{L},\pm}$, following the conventions of~\cite{Sorkin:2019llw,Asante:2021zzh}, where the $\psi_{\sigma,t}$ are dihedral angles. \enquote{$\pm$} indicates a choice of sign which becomes important for causally irregular configurations~\cite{Sorkin:2019llw,Asante:2021zzh,Asante:2021phx} discussed below. The bulk action of Lorentzian Regge calculus is given by~\cite{Asante:2021zzh,Sorkin:2019llw,Regge:1961ct}
%
\begin{equation}\label{eq:general Regge action}
\begin{aligned}
\Sreg[\{l_e\}] =& \sum_{t\in\Delta:\text{ tl}}\abs{A_t(\{l_e\})}\delta_t^{\mathrm{E}}(\{l_e\})\\
+&\sum_{t\in\Delta:\text{ sl}}\abs{A_t(\{l_e\})}\delta_t^{\mathrm{L},\pm}(\{l_e\})\,,
\end{aligned}
\end{equation}
%
with $\abs{A_t}$ the absolute value of the triangle area. In the presence of a boundary, $\partial\Delta\neq\emptyset$, $\Sreg$ attains boundary terms such that overall additivity is ensured. The equations of motion are given by $\pdv{\Sreg}{l_e} = 0$ supplemented by the Schl\"{a}fli identity, $\sum_t \abs{A_t}\pdv{\delta_t}{l_e} = 0$. 

Regge calculus can be formulated in terms of different variables~\cite{Barrett:1994nn,Dittrich:2008va}, with area variables~\cite{Asante:2018wqy,Makela:2000ej,Barrett:1997tx} playing a particularly important role here due to their connection to LQG and spin-foams. The relation between area and length variables is in general not invertible, and area variables need to be further constrained to yield the dynamical equations of length Regge calculus~\cite{Barrett:1997tx,Dittrich:2008va,Asante:2021zzh}. However, in the symmetry reduced setting introduced below, the relation of area and length is globally invertible and the aforementioned constraints trivialize. 

\subsection{Lorentzian 4-frusta}\label{sec:Lorentzian frustum}

Lorentzian $4$-dimensional frusta are ideal to capture spatially flat, homogeneous and isotropic discrete geometries~\cite{Bahr:2017}. In contrast to Chapter~\ref{chapter:specdim} however, we consider here \emph{Lorentzian} $4$-frusta, where the two $3$-cubes are spacelike and where the six boundary $3$-frusta can be either spacelike or timelike. Its entire geometry is captured by the length of the $3$-cube edges, $l_0$ and $l_1$, and its height $H$, defined as the distance between the midpoints of the two $3$-cubes. The general notions of Regge calculus introduced above are straightforwardly adapted to this setting.  

Gluing $4$-frusta along boundary $3$-frusta in spatial direction and along boundary $3$-cubes in temporal direction, one obtains an extended cosmological spacetime of hypercubical combinatorics $\mathcal{X}^{(4)}_\mathcal{V}$ with cubulated slices labelled by $n\in\{0,...,\mathcal{V}\}$. We say that a $4$-frustum lies in a \enquote{slab} between slices $n$ and $n+1$.


The geometry of $4$-frusta as a function of $(l_n,l_{n+1},H_n)$ is detailed in Appendix~\ref{app:Lorentzian frustum}. Most importantly, the squared volume of subcells connecting neighboring slices, i.e. 3-frusta, trapezoids and edges, carries information on the causal structure. If it is positive (negative), the corresponding subcell is timelike (spacelike). In the symmetry reduced setting, the ratio $\frac{l_n-l_{n+1}}{H_n}$ fully characterizes the causal character of subcells suggesting separating the $4$-frusta configurations into three sectors summarized in Fig.~\ref{fig:regions}. Embeddability of $4$-frusta into $\R^{1,3}$ requires the height to be real. Otherwise, the configuration corresponds to a Euclidean $4$-frustum, employed in Chapter~\ref{chapter:specdim}.  As argued below, $H_n$ is naturally associated to the lapse function $N$ which, in the simple context of homogeneous cosmology, relates the Lorentzian and Euclidean sector via an analytical continuation, commonly referred to as Wick rotation.

Dihedral and deficit angles are a crucial ingredient for the Lorentzian Regge action. Here, there are two types of dihedral angles, associated either to spacelike squares or to trapezoids connecting slices. Crucially, the explicit expressions of these angles depend on the causal characters of subcells which are given in Appendix~\ref{app:Lorentzian angles} as a function of $(l_n,l_{n+1},H_n)$.

\subsection{Lorentzian Regge action for spatially flat cosmology}\label{sec:Lorentzian Regge action for spatially flat cosmology}

The underlying combinatorics of the model are hypercubical, denoted by $\mathcal{X}^{(4)}_{\mathcal{V}}$, such that every square or trapezoid is shared by four $4$-frusta. As a result, the Euclidean and Lorentzian exterior boundary deficit angles for a single frustum are defined as $\delta^{\mathrm{E}} = \frac{\pi}{2} - \psi^{\mathrm{E}}$ and $\delta^{\mathrm{L},\pm} = \mp i\frac{\pi}{2} - \psi^{\mathrm{L},\pm}$, respectively, where $\psi$ are the corresponding dihedral angles. Since the single $4$-frustum consists of six squares in slice $n$, twelve trapezoids and six squares in slice $n+1$, spatial homogeneity then implies that the Regge action is given by
%
\begin{equation}\label{eq:slice Regge action}
\begin{aligned}
S_{\mathrm{R}}^{(n)} = 6l_n^2\left(\mp i\frac{\pi}{2}-\varphi_{nn+1}\right)+6l_{n+1}^2\left(\mp i\frac{\pi}{2}-\varphi_{n+1n}\right)+ 12\abs{k_n}\left[\Theta_{\mathrm{sl}}\left(\mp i\frac{\pi}{2}-\theta_n^{\mathrm{L}}\right) + \Theta_{\mathrm{tl}}\left(\frac{\pi}{2} - \theta_n^{\mathrm{E}}\right)\right]\,.
\end{aligned}
\end{equation}
%
Here, $\varphi_{nn+1},\varphi_{n+1n}$ are the dihedral angles located at squares in the $n$th, respectively $(n+1)$th slice and $\theta_n$ is the dihedral angle located at trapezoids. $\Theta_{\mathrm{sl,tl}}=0,1$ is a toggle for the trapezoid being spacelike or timelike.

By construction, the $4$-frustum action is additive such that on a discretization $\mathcal{X}^{(4)}_{\mathcal{V}}$ with $L^3$ vertices in spatial direction and $\mathcal{V}+1$ spatial slices, the action is given by $S_{\mathrm{R}} = L^3\sum_{n=0}^{\mathcal{V}-1}S_{\mathrm{R}}^{(n)}$. Additivity of the action per slab implies that the equations of motion for a given geometric quantity, being either $l_n$ or $H_n$, will only depend on neighboring geometric labels. We discuss the consequences of this slicing structure for the dynamics of the model in more detail in Sec.~\ref{sec:Regge equations and deparametrisation}. 

The Lorentzian Regge action is related to the Euclidean theory discussed previously by a Wick rotation in the height variable $H_n$.  Analogous to continuum studies~\cite{Feldbrugge:2017kzv}, the Wick rotation is performed by complexifying the lapse function, represented in the discrete setting by the squared height $H_n^2\rightarrow R_n^2\e^{i\alpha}$, where $\alpha$ is the rotation angle and $R_n^2$ is the modulus. This procedure has been carried out rigorously first in~\cite{Asante:2021phx}. As a consequence of the Wick rotation, the area of trapezoids as well as the dihedral angles are complexified, denoted by $k_n(\alpha),\theta^\pm_n(\alpha)$ and $\varphi^\pm_{nn+1}(\alpha)$. The resulting Regge action can be analytically continued, the details of which crucially depend on the structure of branch cuts and thus on the causal sector under consideration. Remarkably, the Lorentzian Regge action in Eq.~\eqref{eq:slice Regge action} constructed in an ad hoc manner is directly connected to the Euclidean action derived in~\cite{Bahr:2017eyi} from the semi-classical limit of an underlying quantum geometric model. Although this result suggests that there may exist a symmetry restricted Lorentzian spin-foam model whose semi-classical limit leads to the Lorentzian Regge action of spatially flat cosmology, the ultimate answer to this question remains open, at least in (3+1) dimensions. That is because the semi-classical limit of full Lorentzian spin-foam models, being the EPRL-CH model~\cite{Conrady:2010vx,Conrady:2010kc} or the complete Barrett-Crane model (see Chapter~\ref{chapter:cBC}), is not well understood for timelike interfaces and subject to active research, see~\cite{Simao:2021qno}. In Chapter~\ref{chapter:3d cosmology}, a similar scenario will be studied in (2+1) dimensions where differences of Lorentzian Regge calculus and spin-foam asymptotics can be shown to emerge.

\subsection{Causal (ir)regularity}\label{sec:causal regularity}

The Regge action can attain complex values when the imaginary parts of the dihedral angles at a spacelike face do not sum up to $\mp i2\pi$. This indicates the presence of causal irregularities in the form of a degenerate light cone structure, two examples of which are given by the trouser singularity and the yarmulke singularity associated to spatial topology change~\cite{Horowitz:1990qb,Louko:1995jw,Dowker:1999wu}. Choosing \enquote{$+$} in the convention of Lorentzian angles, the amplitude $\e^{i\Sreg}$ is exponentially suppressed for trouser and exponentially enhanced for yarmulke configurations, vice versa for \enquote{$-$}~\cite{Asante:2021zzh}.  

Causal irregularities can be located at spacelike subcells of different dimension. In (3+1) dimensions, causal regularity conditions are therefore referred to as hinge, edge and vertex causality, respectively~\cite{Asante:2021phx}. In the following, we summarize the results of~\cite{Jercher:2023csk} studying these three types of causality conditions for the present model.

\paragraph{Hinge causality.} Hinge causality violations are characterized by a non-vanishing imaginary part of the Lorentzian deficit angle. Each summand of $\mp i\frac{\pi}{2}$ for a dihedral angle counts one light cone crossing. Thus, there are exactly four light rays at a hinge if the imaginary parts of the dihedral angles sum up to $\mp i2\pi$. At spatial squares, causality violations occur if the 3-frusta are spacelike. At trapezoids, causality violations occur if the trapezoid is spacelike. Overall, the hinge causally regular configurations are found in Sector III (see Fig.~\ref{fig:regions}) where trapezoids and $3$-frusta are timelike. The Regge action in this regime is given by
%
\begin{equation}\label{eq:regular Regge action}
\begin{aligned}
S_{\textsc{iii}}^{(n)} & = 6(l_n^2-l_{n+1}^2)\sinh^{-1}\left(\frac{l_{n+1}-l_n}{\sqrt{4H_n^2-(l_n-l_{n+1})^2}}\right)
& + 12\abs{k_n}\left[\frac{\pi}{2}-\cos^{-1}\left(\frac{(l_n-l_{n+1})^2}{4H_n^2-(l_n-l_{n+1})^2}\right)\right]\,.
\end{aligned}
\end{equation}  
%
Hinge causality is not sensitive to the signature of the edges connecting the $n$th and $(n+1)$th slice, and therefore, both Sector III.1 and III.2 are considered causally regular at the hinges.  

\paragraph{Edge causality.} The causal structure at a spacelike edge $e$ is considered regular if the intersection $P_e\cap\Sigma$ of the orthogonal space of $e$ containing its midpoint and the union $\Sigma = \cup_{\sigma\supset e}\sigma$ of $4$-cells sharing $e$, contains exactly two light cones. The two types of spacelike edges here are 1) those contained in $3$-cubes, and 2) those connecting neighboring slices if $H_n^2<\frac{3}{4}(l_n-l_{n+1})^2$ is met. Following~\cite{Jercher:2023csk} and referring to the details given therein, one finds that the causal structure at edges of type 1) is regular if and only if the trapezoids are timelike, i.e. if $H_n^2 > \frac{1}{2}(l_n-l_{n+1})^2$ holds. On the other hand, the intersection $P_e\cap\Sigma$ for edges of type 2) contains no light cones at all and thus, edge causality is violated everywhere except in Sector III.2.

\paragraph{Vertex causality.} At any vertex $v$ of the cellular complex, consider the union of all $4$-dimensional building blocks that contain that vertex. Then, the causal structure is said to be regular if the intersection of the light cone at $v$ with the boundary of the union are two disconnected spheres~\cite{Asante:2021phx}. Following~\cite{Jercher:2023csk}, this form of causality is also violated everywhere except in Sector III.2.

%Summarizing, Sector I and II exhibit causality violating configurations at vertices, edges and hinges. In Sector III.1 edge and vertex causality are violated which however remains unnoticed by the Regge action, which is only sensitive to hinge causality. 
The results on causality violations in this simplified setting suggest that higher dimensional causality conditions imply lower-dimensional ones, e.g. vertex causality implies hinge causality. It is conceivable that this arises due to the projections onto orthogonal subspaces of $k$-dimensional cells (e.g. edges, $k=1$ or hinges, $k=2$). The relation of different causality violations in a more general setting, however, remains an interesting question for future research.
